\section{Frozen conjectures and boundary distinctions}\label{sec:conj}
The following restates CAND-002 from
\texttt{research/CONJECTURE\_REGISTER.md}, including the
2026-10-27 correction that its third statement is \emph{standalone}.
The original frozen wording and historical clarifications
are preserved in the repository. No statement in this
section is claimed as a theorem.

\begin{conjecture}[CAND-002(1), compact-aspect crossover]\label{conj:one}
There exists \(F:(0,\infty)\to\mathbb R\),
independent of the twist class \(\kappa\), such that for
bounded-factor classes whose exponents all belong to
\(\{1,2,3,4\}\),
\[
 s_\kappa(X)-F\!\left(\frac{\Nm_6(\kappa)}X\right)\longrightarrow0
\]
as \(X\to\infty\), uniformly when \(\Nm_6(\kappa)/X\)
ranges over any fixed compact subset of \((0,\infty)\).
Moreover \(\liminf_{t\to0^+}F(t)>0\), \(F\) is decreasing
on \([10^{-2},1]\), and
\(\limsup_{t\to\infty}F(t)\le0\).
\end{conjecture}

\begin{conjecture}[CAND-002(2), fifth-power large aspect]\label{conj:two}
For \(\kappa=\pi^5\), for every \(\delta>0\)
there exist \(N_0\) and \(T\) such that
\[
 |s_{\pi^5}(X)+1|<\delta
\]
whenever \(\Nm(\pi)\ge N_0\),
\(\Nm(\pi)^5/X\ge T\), and \(X\) is sufficiently large.
\end{conjecture}

\begin{conjecture}[CAND-002(3), standalone fixed classes]\label{conj:three}
For every fixed admissible class \(\kappa\) with
no exponent \(5\),
\(\liminf_{X\to\infty}r_\kappa(X)>1\).
For every fixed prime \(\pi\),
\(\limsup_{X\to\infty}r_{\pi^5}(X)\le1\).
\end{conjecture}

The registry's historical version of (3) phrased it as
a corollary conditional on a sextic Patterson-type
first-moment hypothesis \(H_P\).
The later logical-bridge audit showed that the
\emph{compact-aspect} information in (1) and the
\emph{large-aspect} information in (2) do not control
fixed classes, where \(t\to0\). The 2026-10-27
restatement is the authoritative standalone
fixed-class conjecture. Even corrected \(H_P\) has
errors too large on the typical row scale; see
Section~\ref{sec:heuristic}.
\fig{regime_map}{.90}{The fixed-class boundary is reached by
holding \(\Nm_6\) constant as \(X\) increases, whereas
Conjecture~\ref{conj:one} is uniform on compact positive
\(t\) intervals and Conjecture~\ref{conj:two} concerns
large \(t\). This is a \emph{conceptual} parameter diagram,
not measured data.}{regimes}

Open empirical questions from the frozen record are
(a) whether \(F(0^+)=1\);
(b) the small-aspect fifth-power behavior;
(c) the prime-specific unit-class selection; and
(d) the exponent-two row-parity concentration.
All four are questions, not predictions scored in this paper.
